\documentclass[12pt, a4paper]{article}

% ---- Standard Packages ----
\usepackage[a4paper, margin=1in]{geometry}
\usepackage[utf8]{inputenc}
\usepackage[T1]{fontenc}
\usepackage[bahasa]{babel}
\usepackage{amsmath, amssymb, amsthm, mathtools}
\usepackage{enumitem}
\usepackage{booktabs}
\usepackage{multicol}
\usepackage{hyperref}

% ---- Basic Metadata ----
\title{MODUL BELAJAR BAHASA INDONESIA KELAS 10\\ \large Persiapan Asesmen Sumatif Akhir Tahun (ASAT)}
\author{Ringkasan Materi Komprehensif}
\date{}

\begin{document}

\maketitle

\begin{center}
    \textbf{Topik Pembahasan:} \\
    Teks Laporan Hasil Observasi (LHO) $\cdot$ Teks Anekdot $\cdot$ Puisi \\
    Teks Eksposisi $\cdot$ Afiksasi
\end{center}

\tableofcontents
\newpage

% ==============================================================
\section{TEKS LAPORAN HASIL OBSERVASI (LHO)}
% ==============================================================

\subsection{Pengertian dan Tujuan Teks LHO}
Teks Laporan Hasil Observasi (LHO) adalah teks yang menjabarkan penjabaran umum atau melaporkan sesuatu berupa hasil dari pengamatan (observasi).
\begin{itemize}
    \item \textbf{Tujuan Utama:} Memberikan informasi secara objektif dan faktual mengenai suatu objek, fenomena alam, budaya, atau keadaan sosial kepada pembaca.
    \item \textbf{Objektivitas:} Informasi harus bersumber dari pengamatan nyata, bukan imajinasi atau opini penulis. Kalimat yang mengandung kata-kata seperti \textit{mungkin, sepertinya, menurut saya, cantik, buruk} merupakan kalimat subjektif yang harus dihindari dalam LHO.
\end{itemize}

\subsection{Struktur Teks LHO}
Struktur teks LHO disusun secara sistematis agar mudah dipahami:
\begin{enumerate}
    \item \textbf{Pernyataan Umum (Klasifikasi):} Berisi pengantar, definisi pokok, atau penggolongan objek yang diamati. Contoh kalimat: \textit{"Paus adalah mamalia laut yang termasuk dalam ordo Cetacea."}
    \item \textbf{Deskripsi Bagian:} Berisi penjelasan detail mengenai objek. Jika objeknya hewan, bisa membahas ciri fisik, habitat, atau perilaku. Jika tumbuhan, bisa membahas bentuk daun, akar, bunga, dll.
    \item \textbf{Deskripsi Manfaat:} Menjelaskan fungsi, kegunaan, atau manfaat dari objek yang diamati, baik bagi manusia maupun ekosistem.
\end{enumerate}
\textit{Catatan: Hubungan antarparagraf dalam LHO harus kohesif (terpadu secara bentuk bahasa) dan koheren (terpadu secara makna) melalui penggunaan konjungsi antarparagraf.}

\subsection{Ciri Kebahasaan Teks LHO}
\begin{itemize}
    \item \textbf{Kalimat Definisi:} Menggunakan verba relasional seperti \textit{adalah, merupakan, yaitu, yakni}.
    \item \textbf{Kalimat Deskripsi:} Menggunakan kata-kata yang mendeskripsikan sifat atau wujud fisik objek (contoh: \textit{memiliki kulit bersisik, berukuran 2 meter}).
    \item \textbf{Istilah Teknis (Keilmuan):} Menggunakan kosakata khusus sesuai bidang objek yang diteliti (contoh: \textit{mamalia, ekosistem, abiotik, klorofil}).
    \item \textbf{Kata Benda (Nomina) \& Frasa Nomina:} Banyak menggunakan kata benda sebagai subjek dan objek kalimat.
\end{itemize}

% ==============================================================
\section{TEKS ANEKDOT}
% ==============================================================

\subsection{Pengertian, Tujuan, dan Unsur Humor}
Teks anekdot adalah cerita singkat yang menarik karena kelucuannya, tetapi di balik kelucuan tersebut terdapat pesan berupa kritik atau sindiran terhadap suatu fenomena sosial, tokoh publik, atau instansi.
\begin{itemize}
    \item \textbf{Tujuan Teks:} Menghibur pembaca sekaligus menyampaikan kritik sosial secara halus (tidak menyinggung secara vulgar/kasar).
    \item \textbf{Penyebab Humor:} Kelucuan dalam anekdot biasanya muncul dari elemen kejutan (plot twist), respons tokoh yang absurd, ketidaknyambungan komunikasi, atau permainan kata (pelesetan).
    \item \textbf{Kritik dan Sindiran:} Sasaran sindiran umumnya adalah ketidakadilan, korupsi, buruknya layanan publik, atau perilaku tokoh berkuasa. Efektivitas anekdot dinilai dari kemampuan pesan kritik tersebut tersampaikan melalui medium humor.
\end{itemize}

\subsection{Struktur Teks Anekdot}
\begin{enumerate}
    \item \textbf{Abstraksi:} Bagian awal pendahuluan yang menyatakan latar belakang atau gambaran umum cerita.
    \item \textbf{Orientasi:} Bagian yang mengarah pada terjadinya suatu peristiwa atau konflik awal.
    \item \textbf{Krisis/Komplikasi:} Inti peristiwa; tempat terjadinya kekonyolan, masalah, atau kejanggalan yang memuat kritik atau sindiran.
    \item \textbf{Reaksi:} Langkah penyelesaian atau tanggapan dari tokoh terhadap masalah pada bagian krisis (biasanya berupa respons yang memancing tawa).
    \item \textbf{Koda:} Penutup atau kesimpulan cerita (bisa ada, bisa tidak).
\end{enumerate}

\subsection{Ciri Kebahasaan Teks Anekdot}
\begin{itemize}
    \item Menggunakan keterangan waktu lampau (contoh: \textit{Pada suatu hari, suatu ketika}).
    \item Menggunakan pertanyaan retoris (pertanyaan yang tidak membutuhkan jawaban).
    \item Menggunakan konjungsi temporal (waktu) seperti \textit{kemudian, lalu, setelah itu}.
    \item Menggunakan kata kerja material (aksi) seperti \textit{berjalan, menulis, terdiam}.
    \item Menggunakan kalimat seru, kalimat perintah, dan majas ironi/sarkasme.
\end{itemize}

% ==============================================================
\section{PUISI}
% ==============================================================

\subsection{Karakteristik dan Jenis Puisi Lama}
Puisi lama sangat terikat oleh aturan jumlah baris, rima (sajak), suku kata, dan irama.
\begin{itemize}
    \item \textbf{Pantun:} 4 baris/bait, sajak a-b-a-b, baris 1-2 sampiran, 3-4 isi. Mengandung nasihat, jenaka, atau teka-teki.
    \item \textbf{Syair:} 4 baris/bait, sajak a-a-a-a, semua baris adalah isi. Cenderung mengisahkan cerita sejarah, agama, atau dongeng.
    \item \textbf{Gurindam:} 2 baris/bait, sajak a-a. Baris pertama syarat/masalah, baris kedua akibat/jawaban. Sangat sarat akan petuah hidup dan nilai agama.
    \item \textbf{Karmina (Pantun Kilat):} Hanya 2 baris (baris 1 sampiran, baris 2 isi), rima a-a.
\end{itemize}

\subsection{Unsur Pembangun Puisi (Makna, Diksi, Suasana, Amanat)}
\begin{itemize}
    \item \textbf{Diksi \& Makna Simbolik:} Penyair memilih kata-kata kiasan (konotatif) untuk memperdalam makna. Makna simbolik adalah lambang tertentu, misalnya \textit{"fajar"} menyimbolkan \textit{harapan baru}, atau \textit{"lintah darat"} menyimbolkan \textit{rentenir}.
    \item \textbf{Suasana Puisi:} Keadaan psikologis atau emosi yang dirasakan pembaca (misal: syahdu, sedih, bersemangat, marah, haru).
    \item \textbf{Kritik Sosial \& Amanat:} Puisi sering kali menjadi medium protes terhadap ketidakadilan (kritik sosial). Amanat adalah pesan moral dari penyair kepada pembaca.
\end{itemize}

\subsection{Pembacaan dan Musikalisasi Puisi}
\begin{itemize}
    \item \textbf{Aspek Pembacaan (Deklamasi):} Penilaian berfokus pada \textit{vokal} (artikulasi, intonasi, volume, jeda), \textit{ekspresi} (penghayatan, mimik wajah), dan \textit{kinesik} (gerak tubuh).
    \item \textbf{Musikalisasi Puisi:} Pembacaan puisi yang dipadukan dengan alunan musik. Musik berfungsi untuk mempertegas suasana puisi, bukan mendominasi vokalnya.
\end{itemize}

% ==============================================================
\section{TEKS EKSPOSISI}
% ==============================================================

\subsection{Pengertian dan Tujuan Teks Eksposisi}
Teks eksposisi adalah teks yang bertujuan untuk menjelaskan, menguraikan, atau memaparkan suatu gagasan dan informasi secara jelas agar pengetahuan pembaca bertambah. Teks ini bersifat logis dan faktual.

\subsection{Struktur Teks Eksposisi}
\begin{enumerate}
    \item \textbf{Tesis (Pernyataan Pendapat):} Gagasan utama, pandangan, atau prediksi penulis tentang suatu permasalahan berdasarkan fakta awal.
    \item \textbf{Argumentasi:} Alasan-alasan logis berupa fakta, data statistik, atau pendapat ahli yang mendukung dan memperkuat tesis.
    \item \textbf{Penegasan Ulang (Reiterasi):} Kesimpulan yang merangkum keseluruhan argumen dan menegaskan kembali tesis tanpa menambahkan argumen baru. Sering disertai saran atau rekomendasi.
\end{enumerate}

\subsection{Fakta vs Opini dalam Teks Eksposisi}
\begin{itemize}
    \item \textbf{Fakta:} Pernyataan yang benar-benar terjadi, memiliki bukti/data kuantitatif, atau berupa peristiwa masa lalu. \textit{Contoh: "Berdasarkan data BPS tahun 2023, angka kemiskinan turun sebesar 1,2\%."}
    \item \textbf{Opini:} Gagasan, pemikiran, penilaian, atau harapan penulis. Biasanya ditandai dengan kata: \textit{sebaiknya, seharusnya, mungkin, jika, menurut saya}. \textit{Contoh: "Pemerintah seharusnya lebih tegas dalam menangani kasus korupsi."}
\end{itemize}

% ==============================================================
\section{AFIKSASI (PENGIMBUHAN)}
% ==============================================================

\subsection{Konsep Dasar Afiksasi}
Afiksasi adalah proses morfologi berupa penambahan imbuhan (afiks) pada bentuk dasar sehingga menghasilkan kata baru (kata bentukan/turunan) yang mengalami perubahan kelas kata atau makna.

\subsection{Proses dan Jenis Afiksasi}
\begin{itemize}
    \item \textbf{Prefiks (Awalan):} Imbuhan di awal kata.
        \begin{itemize}
            \item Contoh: \textit{me-} (menulis), \textit{di-} (dibaca), \textit{ber-} (berlari), \textit{pe-} (penulis).
        \end{itemize}
    \item \textbf{Sufiks (Akhiran):} Imbuhan di akhir kata.
        \begin{itemize}
            \item Contoh: \textit{-an} (tulisan, makanan), \textit{-kan} (bacakan), \textit{-i} (turuti).
        \end{itemize}
    \item \textbf{Infiks (Sisipan):} Imbuhan yang disisipkan di tengah kata dasar (setelah konsonan pertama). Dalam bahasa Indonesia, infiks meliputi \textbf{-el-}, \textbf{-em-}, \textbf{-er-}, dan \textbf{-in-}.
        \begin{itemize}
            \item Contoh infiks \textit{-em-}: \textit{tali} $\rightarrow$ t\textbf{em}ali, \textit{kilau} $\rightarrow$ k\textbf{em}ilau.
            \item Contoh infiks \textit{-el-}: \textit{tapak} $\rightarrow$ t\textbf{el}apak.
            \item Contoh infiks \textit{-er-}: \textit{suling} $\rightarrow$ s\textbf{er}uling, \textit{gigi} $\rightarrow$ g\textbf{er}igi.
            \item Contoh infiks \textit{-in-}: \textit{kerja} $\rightarrow$ k\textbf{in}erja.
        \end{itemize}
    \item \textbf{Konfiks (Gabungan Awalan dan Akhiran Padu):} Imbuhan yang ditambahkan di awal dan akhir kata dasar \textbf{secara serentak} sebagai satu kesatuan. Jika salah satu dilepas, kata tersebut seringkali tidak memiliki makna baku.
        \begin{itemize}
            \item Contoh: \textit{ke-an} pada kata \textit{keadilan} (dari \textit{adil}).
            \item Contoh: \textit{pe-an} pada kata \textit{pendaftaran} (dari \textit{daftar}).
            \item Contoh: \textit{per-an} pada kata \textit{perdamaian} (dari \textit{damai}).
        \end{itemize}
\end{itemize}

\end{document}
